\section{Institutional Setting and Data}\label{sec:data}

\subsection{Scheduled releases and the consensus}

The U.S.\ statistical agencies and private survey providers publish their main indicators on
calendars fixed well in advance. Most price and labor-market data, including the Consumer Price
Index (CPI) and the Employment Situation report (nonfarm payrolls, NFP), are released at 08:30
Eastern Time; survey-based indicators such as the ISM indices and consumer confidence are released
at 10:00; and FOMC rate decisions at 14:00. Before each release, Bloomberg surveys professional
forecasters and publishes the median forecast, which we use as the market consensus. Because the
consensus is fixed before the release and the actual figure is published at a known second, the
difference between the two is news to the market at the moment of the release.

\subsection{Futures data}

We use one-minute OHLCV bars from Databento (CME Globex, \texttt{GLBX.MDP3}) for four contracts
from July 2010 to June 2026: the E-mini S\&P 500 (\ES), E-mini Nasdaq-100 (\NQ) and E-mini Dow (\YM),
rolled on the calendar, and the 10-year Treasury note future (\ZN), rolled on volume. Table~\ref{tab:contracts}
summarizes the contract specifications. The markets differ in how they should react to the
same news. \YM, the Dow, holds fewer technology and more industrial and financial stocks than
\ES, and serves as a third equity market on which the \ES\ results can be checked. A hot inflation print raises the expected path of the policy rate, which lowers the
present value of future earnings---more so for the long-duration technology stocks that dominate
\NQ---and lowers Treasury prices. Strong payrolls raise expected growth, which is good for
equities, but also expected rates, which is bad for bonds.

\begin{table}[htbp]
  \centering\small
  \caption{Contract specifications and round-trip costs.}\label{tab:contracts}
  \begin{threeparttable}
  \begin{tabular}{lcccc}
    \toprule
    & \ES & \NQ & \YM & \ZN \\
    \midrule
    Underlying & S\&P 500 & Nasdaq-100 & Dow Jones Ind. & 10-year note \\
    Exchange & CME & CME & CBOT & CBOT \\
    Tick size & 0.25 points & 0.25 points & 1 point & 1/64 point \\
    Value of one tick & \$12.50 & \$5.00 & \$5.00 & \$15.625 \\
    Round-trip cost, 2 ticks + \$4.50 (bps) & $\approx 1.7$ & $\approx 0.7$ & $\approx 1.0$ & $\approx 3.0$ \\
    Buy-and-hold Sharpe ratio, 2019--2026 & 0.90 & 0.99 & 0.78 & $-0.23$ \\
    Buy-and-hold maximum drawdown, 2019--2026 & $-29.5\%$ & $-38.3\%$ & $-32.2\%$ & $-23.6\%$ \\
    \bottomrule
  \end{tabular}
  \begin{tablenotes}\footnotesize
    \item Notes: The round-trip cost converts two ticks of slippage (one per side) plus a \$4.50
    commission into basis points of the contract value at the pre-release price, averaged over
    2019--2026.
  \end{tablenotes}
  \end{threeparttable}
\end{table}

We keep sessions from 18:00 to 17:00 Eastern Time, dated by their close, and drop exchange
holidays and sessions with too few bars. Minutes without a trade inside a kept session are filled
with the last price. Kept sessions number 3,894 for \ES, 3,797 for \NQ, 3,754 for \YM\ and 3,916
for \ZN; together the four markets contribute about 21 million one-minute bars. Returns across a change of contract are set to zero, which removes the
calendar-spread jump at each quarterly roll and leaves returns within a contract unchanged. A
one-minute bar is labelled by its start: the bar labelled 08:30 covers 08:30:00--08:30:59 and
closes at 08:31. Figure~\ref{fig:markets} shows the four markets over the sample, which spans the
low-inflation years before 2020, the COVID crash, the 2022 inflation shock with the fastest
monetary tightening in four decades, and the disinflation that followed.

\begin{figure}[htbp]
  \centering
  \includegraphics[width=0.85\linewidth]{fig01_markets_overview_4m}
  \caption{\ES, \NQ, \YM\ and \ZN, 2010--2026. Roll-adjusted daily prices rebased to 100. Grey: the
  2019--2026 evaluation period; red: the 2022 inflation shock.}
  \label{fig:markets}
\end{figure}

\subsection{Measuring surprises}

The Bloomberg calendar contains 179 U.S.\ release lines and 27,137 release rows between 2010 and
2026, each with a release time, the actual value and the survey median. For line $l$ at release
$t$, the raw surprise $A_{l,t}-F_{l,t}$ is not comparable across lines (payrolls are measured in
thousands of jobs, CPI in tenths of a percent), so we standardize it causally, using only earlier
releases of the same line:
\begin{equation}\label{eq:z}
  z_{l,t} = \frac{A_{l,t}-F_{l,t}}{\operatorname{sd}\left(A_{l,1}-F_{l,1},\ldots,A_{l,t-1}-F_{l,t-1}\right)},
\end{equation}
requiring at least 24 earlier releases and capping $z$ at $\pm 5$ (0.47\% of the 17,328
standardized surprises are capped). Lines that are released at the same instant---for example the
eight CPI lines---are grouped into a release family $f$ when their release timestamps overlap by at
least 80\%. The family signal is
\begin{equation}\label{eq:x}
  x_{f,t} = \frac{1}{|L_f|}\sum_{l\in L_f} s_l\, z_{l,t},
\end{equation}
where the sign $s_l\in\{-1,+1\}$ is the sign of the slope of a regression of the first-minute
return on $z_{l,t}$, estimated separately for each market on earlier years only. With these signs,
$x_{f,t}>0$ always means good news for the market in question: a strong payrolls print is good
news for \ES\ and bad news for \ZN.

\subsection{Sample}

The event-study sample (Section~\ref{sec:discovery}) contains 363 CPI, NFP and PCE releases from
2016 to 2026. The trading samples use all release families with at least 30 releases (75 families
built from 143 lines). All strategies use the sample split in Table~\ref{tab:split}; the first
three years only build the history that later estimates require.

\begin{table}[htbp]
  \centering\small
  \caption{Sample split.}\label{tab:split}
  \begin{tabular}{lll}
    \toprule
    Period & Years & Use \\
    \midrule
    Burn-in    & 2010--2012 & analogue library and surprise history only \\
    Train      & 2013--2017 & first estimation window \\
    Validation & 2018--2020 & out of sample (walk-forward) \\
    Test       & 2021--2026 & out of sample, reported once \\
    \bottomrule
  \end{tabular}
\end{table}
